\chapter*{ABSTRACT}
\addcontentsline{toc}{chapter}{Abstract}

The recruitment industry faces significant structural challenges in processing the large volume of unstructured applications submitted for modern job postings. Online job portals, professional networking platforms, and remote hiring channels have increased the screening burden on recruiters. Manual resume screening is time-consuming, resource-intensive, and vulnerable to inconsistency and unconscious bias. Under time constraints, recruiters may depend on visual formatting, recognizable institutions, or familiar company names rather than a systematic analysis of candidate-job fit.

Automated Applicant Tracking Systems (ATS) reduce this burden, but many legacy systems rely primarily on rigid keyword matching. Such systems often disregard semantic context and fail to capture implicit candidate competencies. Conversely, Large Language Models (LLMs) and sentence embeddings such as Sentence-BERT provide stronger semantic understanding but are computationally expensive, memory-intensive, slower for bulk processing, and often operate as ``black boxes'' whose scores are difficult to justify or audit.

This project explores the strategic implementation of Artificial Intelligence within resume analysis software to balance matching accuracy, computational efficiency, and interpretability. It does not claim to produce a complete commercial hiring replacement or a universally superior AI model. Instead, it develops a lightweight, transparent, and cost-efficient prototype for first-pass screening and decision support in constrained hardware environments and small-to-medium enterprise (SME) workflows.

The proposed pipeline combines multiple lightweight AI and software engineering components. Resumes and job descriptions are accepted through a Flask-based web interface. The parser supports digital PDF, DOCX, TXT, and image-based resume formats. To reduce format bias, the system uses digital parsing where possible and triggers an OCR fallback pipeline for scanned PDFs or image resumes. This pipeline applies Pillow-based grayscale conversion, binarization, and sharpening before extracting text using Tesseract OCR.

After text extraction, Natural Language Processing (NLP) is performed using SpaCy and rule-based regular expression matching. The system extracts technical skills, soft skills, and years of experience from both the job description and candidate resumes. For contextual similarity, the project uses Term Frequency-Inverse Document Frequency (TF-IDF) vectorization with Cosine Similarity. A custom weighted scoring formula combines skill match, experience match, contextual similarity, and an extra skills bonus. The score is bounded so that overqualified experience or excessive additional skills do not distort the result.

To bridge the algorithmic transparency gap, a custom Explainable Artificial Intelligence (XAI) module translates numerical matching vectors into human-readable score breakdowns. Instead of returning only a final percentage, the system lists matched skills, missing skills, extra skills, experience alignment, and score component contributions. This positions the prototype as a recruiter-assistive tool rather than an autonomous selection system.

Benchmarking across five experimental setups evaluated the prototype and compared architectural trade-offs. Benchmark 1 tested format bias and OCR failover resilience across TXT, PDF, and PNG inputs. Benchmark 2 evaluated resistance to dilution and keyword stuffing. Benchmark 3 validated XAI output and mathematical boundary conditions. Benchmark 4 compared the implemented Regex + TF-IDF prototype (Approach A), a pure Sentence-BERT/LLM semantic approach (Approach C), and a future hybrid approach combining regex-based hard boundaries with LLM-based contextual scoring (Approach B). Benchmark 5 measured efficiency, scalability, and cost across 100 resumes.

The quantitative findings reveal a clear trade-off. LLM-based models perform better at resolving semantic gaps, such as understanding that ``RESTful'' and ``REST'' may represent related concepts. However, the implemented lexical prototype operates in milliseconds, achieving a measured processing time of 0.0062 seconds per resume with a near-zero memory footprint of 0.12 MB. In comparison, the pure LLM approach required 4.1860 seconds and 7.33 MB peak RAM, while the future hybrid approach required 3.7217 seconds and 1.35 MB peak RAM. Thus, although the prototype is not semantically superior to transformer-based systems, it is more efficient, transparent, and suitable for low-cost first-pass filtering.

Ultimately, this project validates a transparent decision-support system that accelerates initial shortlisting while maintaining human-in-the-loop oversight. It demonstrates that recruitment automation does not always require the heaviest computational model; it requires an engineering approach aligned with deployment context, resource constraints, transparency needs, and user risk.

\clearpage
